\section{The standard \v{C}ech computation on \texorpdfstring{$\mathbf P^n$}{Pn}}

\subsection{The finite complex}

Let $n\geq1$ and let
\[
           \mathbf P^n_R=\bigcup_{i=0}^{n}U_i,
           \qquad U_i=D_+(X_i).
\]
Every finite intersection $U_{i_0}\cap\cdots\cap U_{i_p}$ is a standard affine open.  In homogeneous
coordinates, its degree-$d$ sections of $\cO(d)$ are finite Laurent polynomials in which
$X_{i_0},\ldots,X_{i_p}$ may have negative exponents and all other variables have nonnegative
exponents, with total degree $d$.

The ordered \v{C}ech complex is
\begin{equation}\label{eq:pn-cech}
 \check C^p_d=\bigoplus_{0\leq i_0<\cdots<i_p\leq n}
 \left(R[X_0,\ldots,X_n,(X_{i_0}\cdots X_{i_p})^{-1}]\right)_d,
\end{equation}
with the alternating restriction differential.

Every individual code in \eqref{eq:pn-cech} has finite Laurent support.  Thus an iPRA theorem does
not form an actually infinite direct sum: it scans the finitely many exponent vectors occurring in
the input cochain.

\subsection{Decomposition by exponent vector}

Fix $e=(e_0,\ldots,e_n)\in\Z^{n+1}$ with
\[
                            \sum_{i=0}^{n}e_i=d.
\]
Let
\[
                     \operatorname{Neg}(e)=\set{i:e_i<0}.
\]
The monomial $X^e=X_0^{e_0}\cdots X_n^{e_n}$ occurs in the component indexed by
$I=\{i_0<\cdots<i_p\}$ exactly when
\[
                       \operatorname{Neg}(e)\subseteq I.
\]
Consequently the full complex is the coefficientwise direct sum of finite complexes
$C^\bullet(e)$, each having one copy of $R$ for every subset $I$ that contains
$\operatorname{Neg}(e)$.

The differential on $C^\bullet(e)$ is purely combinatorial and independent of the coefficients:
\[
       (ds)_J=\sum_{a=0}^{|J|-1}(-1)^a s_{J\setminus\{j_a\}},
\tag{12.1}
\]
where a missing component is interpreted as zero.

\subsection{The three sign cases}

\begin{theorem}[Monomial summands]\label{thm:monomial-summands}
For a fixed exponent vector $e$:
\begin{enumerate}[label=\arabic*.]
\item If every $e_i\geq0$, then $C^\bullet(e)$ has cohomology $R$ in degree $0$ and zero in positive
      degrees.
\item If every $e_i<0$, then $C^\bullet(e)$ has a single nonzero term, in degree $n$, and its
      cohomology is $R$ there.
\item If $\operatorname{Neg}(e)$ is nonempty and proper, then $C^\bullet(e)$ is contractible by an
      explicit primitive-recursive homotopy.
\end{enumerate}
\end{theorem}

\begin{proof}
In the first case, every subset component is present.  This is the ordered cochain complex of a full
$n$-simplex.  Its degree-zero constant tuples give one copy of $R$, and the usual insertion homotopy
contracts the positive degrees.

In the second case, a subset containing all $n+1$ indices must be the full set, so the only term is in
degree $n$.

For the mixed case, choose the least index
$i_*\notin\operatorname{Neg}(e)$.  For a $p$-cochain $s$, define the $(p-1)$-cochain $hs$ on an
ordered set $I=\{i_0<\cdots<i_{p-1}\}$ by zero when $i_*\in I$, and otherwise by inserting $i_*$:
\begin{equation}\label{eq:insertion-homotopy}
 (hs)_I=(-1)^a s_{I\cup\{i_*\}},
 \qquad
 a=\#\set{i\in I:i<i_*}.
\end{equation}
The component on the right exists because adding $i_*$ does not change whether the set contains
$\operatorname{Neg}(e)$.  A finite sign computation pairs every deletion not involving $i_*$ once in
$dh$ and once with the opposite sign in $hd$; the remaining deletion of $i_*$ contributes $s_I$.
Hence
\[
                              dh+hd=\mathrm{id}.
\]
All sums have at most $n+1$ terms.
\end{proof}

\subsection{Cohomology table}

\begin{theorem}[Finite \v{C}ech cohomology of projective space]\label{thm:pn-cohomology}
For $n\geq1$, the standard-cover \v{C}ech cohomology is
\[
 \check H^q(\mathbf P^n_R,\cO(d))\cong
 \begin{cases}
   R[X_0,\ldots,X_n]_d,& q=0,\ d\geq0,\\[0.2em]
   \displaystyle
   \left(\frac1{X_0\cdots X_n}
   R[X_0^{-1},\ldots,X_n^{-1}]\right)_d,& q=n,\ d\leq-n-1,\\[0.5em]
   0,&\text{otherwise.}
 \end{cases}
\]
The two nonzero modules are finite free.  Their ranks are
\[
    \operatorname{rank}H^0=\binom{n+d}{n}\quad(d\geq0),
    \qquad
    \operatorname{rank}H^n=\binom{-d-1}{n}\quad(d\leq-n-1).
\]
The isomorphisms and zero contractions are primitive-recursive on finite cochain codes.
\end{theorem}

\begin{proof}
Decompose every finite cochain coefficientwise by exponent vector.  By
Theorem~\ref{thm:monomial-summands}, only all-nonnegative exponent vectors survive in degree zero and
all-negative vectors survive in degree $n$.  The all-nonnegative vectors of total degree $d$ are the
ordinary degree-$d$ monomials, which exist only for $d\geq0$ and are counted by
$\binom{n+d}{n}$.

For an all-negative vector write $e_i=-a_i-1$ with $a_i\geq0$.  Then
\[
                   \sum_i a_i=-d-(n+1).
\]
Such vectors exist exactly when $d\leq-n-1$, and stars-and-bars counts them by
$\binom{-d-1}{n}$.  The normal-form map keeps precisely those monomial coefficients and applies the
insertion homotopy to every mixed-sign summand.
\end{proof}

For $n=1$, Theorem~\ref{thm:pn-cohomology} specializes to
Theorems~\ref{thm:p1-h0} and~\ref{thm:p1-h1}.

\subsection{Primitive-recursive realization}

The implementation needs the following fixed operations:

\begin{itemize}
\item signed exponent-vector coding and bounded coordinate access;
\item the set $\operatorname{Neg}(e)$ as a bit mask;
\item enumeration of subsets of $\{0,\ldots,n\}$ by bit masks;
\item alternating-sign finite sums;
\item insertion of the least index outside a bit mask;
\item finite maps from exponent vectors to coefficient tuples;
\item coefficientwise splitting into all-nonnegative, all-negative, and mixed classes.
\end{itemize}

All are primitive recursive.  To prove \eqref{eq:insertion-homotopy}, an object-iPRA development can
avoid abstract matrices.  For each output component and each deletion index, a bounded case split
checks whether the deleted index is $i_*$, below it, or above it.  A finite accumulator records the
paired terms and their signs.  The induction formula is an equality of natural-number-coded finite
sums.

\subsection{Executable matrix check}

The companion program independently constructs the finite matrices of $d$ and $h$ for every proper
nonempty negative-index set in dimensions $1\leq n\leq6$.  It verifies
$dh+hd=I$ over the integers in 240 mixed-sign cases.  This is not a proof for all $n$, but it checks
the exact sign convention and the finite-index implementation used by the Lean plan.

\subsection{Conventional cohomology}

The standard cover has affine finite intersections, and $\cO(d)$ is quasi-coherent.  The conventional
comparison theorem therefore identifies the finite \v{C}ech calculation with
$H^q(\mathbf P^n_R,\cO(d))$.  The Stacks Project gives the same exponent-vector decomposition and
insertion homotopy.  The new foundational point is that the entire calculation itself is finite,
explicit, and admits a direct route to quantifier-free induction.
